% =========================================================================
% ปฏิบัติการที่ 6: การเชื่อมต่อระบบปัญญาประดิษฐ์และเซนเซอร์ IoT
% =========================================================================

\chapter{การเชื่อมต่อระบบปัญญาประดิษฐ์และเซนเซอร์ IoT}
\label{ch:lab06}

\noindent{\large\bfseries\color{cyberblue} การรับส่งสตรีมข้อมูลเซนเซอร์ผ่าน WebSocket และเรนเดอร์แดชบอร์ด 3D Spatial Hologram}

\vspace{0.4cm}

\begin{labobjectivebox}{การเชื่อมต่อระบบปัญญาประดิษฐ์และเซนเซอร์ IoT}
\textbf{วัตถุประสงค์การเรียนรู้}
\begin{enumerate}[leftmargin=1.8cm, itemsep=2pt]
    \item เข้าใจการรับส่งข้อมูลเรียลไทม์ระหว่างอุปกรณ์ IoT และ WebAR ผ่านโปรโตคอล WebSocket / MQTT
    \item สามารถแปลงข้อมูลอนุกรมเวลา (Time-series Sensor Data) เป็นแผนภูมิ 3 มิติเชิงพื้นที่
    \item สามารถสร้างระบบดิจิทัลทวิน (Digital Twin) พื้นฐานที่ตอบสนองต่อสภาวะแวดล้อมจริง
\end{enumerate}

\vspace{4pt}
\textbf{เครื่องมือและอุปกรณ์จำเป็น} \enspace ไมโครคอนโทรลเลอร์ ESP32 พร้อมเซนเซอร์ DHT22/BME280 หรือจำลองผ่าน WebSocket Server
\end{labobjectivebox}

\section{หลักการและทฤษฎีพื้นฐาน}

การบูรณาการระบบกายภาพและระบบไซเบอร์เข้าด้วยกัน (Cyber-Physical Systems: CPS) เป็นหัวใจของเทคโนโลยีดิจิทัลทวิน (Digital Twin) ในการทดลองนี้ ข้อมูลกายภาพจากเซนเซอร์ IoT เช่น ค่าอุณหภูมิ ความชื้น และค่าฝุ่นละอองขนาดเล็ก จะถูกส่งเข้าสู่เว็บแอปพลิเคชันผ่านช่องทางสื่อสารสองทิศทางความเร็วสูง (WebSocket Protocol):
\begin{equation}
\text{Data Stream}: \{ t_k, T_k, H_k, \text{PM}_{2.5}^{(k)} \} \xrightarrow{\text{WebSocket}} \text{WebAR Scene}
\end{equation}
ในฝั่ง WebAR ข้อมูลที่ได้รับจะถูกนำไปอัปเดตสถานะของแบบจำลอง 3 มิติ เช่น การเปลี่ยนสีของการ์ดแสดงผลโฮโลกราฟิก การปรับเปลี่ยนความหนาแน่นของอนุภาคละอองลอย (Particle System) เพื่อจำลองมลพิษในอากาศตามค่าจริง และการส่งสัญญาณเตือนภัยด้วยแสงและเสียงเมื่อตรวจพบค่าที่ผิดปกติผ่านการวิเคราะห์ด้วยโมเดลการเรียนรู้ของเครื่อง

\begin{conceptbox}{สาระสำคัญของปฏิบัติการที่ 6}
การเข้าใจโครงสร้างทางคณิตศาสตร์และการประมวลผลเชิงพื้นที่ ช่วยให้นักศึกษาสามารถออกแบบระบบเสมือนจริงที่ตอบสนองต่อผู้ใช้งานได้อย่างเป็นธรรมชาติและแม่นยำสูง ทั้งยังสามารถต่อยอดสู่การพัฒนาแอปพลิเคชันทางวิทยาศาสตร์และอุตสาหกรรมได้อย่างมีประสิทธิภาพ
\end{conceptbox}

\section{ขั้นตอนการปฏิบัติการและการทดลอง}

ให้นักศึกษาดำเนินการสร้างไฟล์โครงการและเขียนคำสั่งตามลำดับขั้นตอนต่อไปนี้ โดยตรวจสอบความถูกต้องของไลบรารีที่เรียกใช้งานและเส้นทางไฟล์

\begin{codebox}{โครงสร้างโค้ดหลักของปฏิบัติการที่ 6}
\begin{lstlisting}[language=HTML]
// การเชื่อมต่อ WebSocket Client ใน WebAR
const socket = new WebSocket('ws://localhost:8080');

socket.onopen = () => {
    console.log('เชื่อมต่อเซิร์ฟเวอร์ IoT เรียบร้อย');
};

socket.onmessage = (event) => {
    const telemetry = JSON.parse(event.data);
    updateSpatialDashboard(telemetry);
};

function updateSpatialDashboard(data) {
    // ปรับเปลี่ยนค่าพารามิเตอร์ของวัตถุ 3 มิติตามข้อมูลเซนเซอร์
    const temp = data.temperature;
    const humidity = data.humidity;
    
    // ปรับสีวัตถุตามอุณหภูมิ: เย็น (สีฟ้า) -> ร้อน (สีแดง)
    const normalizedTemp = Math.min(Math.max((temp - 20) / 20, 0), 1);
    const targetColor = new THREE.Color().lerpColors(
        new THREE.Color(0x06b6d4), // ฟ้า
        new THREE.Color(0xef4444), // แดง
        normalizedTemp
    );
    cube.material.color = targetColor;
    
    // อัปเดตข้อความ 3D Spatial Canvas
    updateTextTexture(`Temp: ${temp} C | Humid: ${humidity}%`);
}
\end{lstlisting}
\end{codebox}

\section{การทดสอบระบบและการบันทึกผล}

เมื่อรันโปรแกรมบนเซิร์ฟเวอร์จำลอง (Local Development Server) ให้ทำการทดสอบการทำงานตามเกณฑ์ประเมินในตารางต่อไปนี้ พร้อมบันทึกผลการสังเกตการณ์ลงในรายงานผลการทดลอง

\begin{table}[htbp]
\centering
\small
\begin{tabularx}{\textwidth}{c X c Y}
\toprule
\textbf{ลำดับ} & \textbf{รายการทดสอบและพฤติกรรมที่คาดหวัง} & \textbf{เกณฑ์ผ่าน} & \textbf{ผลการทดสอบ} \\
\midrule
1 & การเริ่มต้นระบบและการเข้าถึงสิทธิ์ฮาร์ดแวร์ & ภายใน 3 วินาที & ผ่าน / ไม่ผ่าน \\
2 & ความเสถียรของการตรวจจับและการคำนวณพิกัด & อัตราความแม่นยำ $> 90\%$ & ผ่าน / ไม่ผ่าน \\
3 & อัตราการแสดงผลกราฟิกต่อเนื่อง (Framerate) & ไม่ต่ำกว่า 55 FPS & ผ่าน / ไม่ผ่าน \\
4 & การตอบสนองต่อปฏิสัมพันธ์เชิงพื้นที่ & ความหน่วง $< 40$ ms & ผ่าน / ไม่ผ่าน \\
\bottomrule
\end{tabularx}
\caption{ตารางประเมินผลการทำงานของระบบในปฏิบัติการที่ 6}
\label{tab:eval_lab06}
\end{table}

\section{ภารกิจท้าทายและการต่อยอด}

\begin{activitybox}{กิจกรรมส่งเสริมทักษะขั้นสูง}
ให้นักศึกษาเขียนโค้ดเพิ่มระบบจำลองอนุภาคฝุ่นควัน (Smoke Particle System) ใน Three.js โดยให้ความเร็วในการลอยตัวและจำนวนอนุภาคแปรผันตามค่าตัวเลขมลพิษที่ได้รับจากเซนเซอร์
\end{activitybox}

\section{คำถามท้ายการทดลอง}

ให้นักศึกษาตอบคำถามเชิงวิเคราะห์ต่อไปนี้ลงในสมุดบันทึกปฏิบัติการ

\begin{enumerate}[leftmargin=1.8cm, itemsep=6pt]
    \item เปรียบเทียบข้อดีและข้อเสียของการใช้ WebSocket เทียบกับ HTTP REST Polling ในงานแสดงผลค่าเซนเซอร์เชิงพื้นที่
    \item แนวคิดของ ดิจิทัลทวิน (Digital Twin) แตกต่างจากการทำกราฟแสดงผลแดชบอร์ด 2 มิติทั่วไปอย่างไร
    \item หากการเชื่อมต่อเครือข่ายขาดหายไประหว่างการแสดงผล ควรมีกลยุทธ์การคืนสภาพ (Failover \& State Reconnect) ใน WebAR อย่างไร
\end{enumerate}

